\chapter{Creio em Deus Pai todo-poderoso, criador do céu e da terra}
\chaptermark{Creio em Deus Pai todo-poderoso}

Este capítulo aprofunda o primeiro artigo do Símbolo Apostólico, professando a fé em Deus como Pai todo-poderoso e Criador de tudo o que existe. Desdobramos os parágrafos iniciais do Catecismo da Igreja Católica (CIC), revelando a natureza de Deus, sua obra criadora, os anjos, a dignidade humana e as origens do mal. Para adultos em processo de iniciação cristã, essa meditação fortalece a conversão, preparando o coração para o Batismo, Crisma e Eucaristia na Vigília Pascal.

\section{A Profissão de Fé no Deus Único (\S 232--267)}

A profissão de fé cristã inicia com as palavras ``Creio em Deus'', um ato pessoal e comunitário que expressa a adesão total do ser humano ao Mistério revelado. Essa declaração não é uma simples opinião intelectual, mas uma resposta confiante à autocomunicação de Deus na história, enraizada na Escritura e na Tradição apostólica.

O ato de fé envolve a inteligência, que acolhe a verdade divina, e a vontade, que se entrega livremente a Deus. É um dom do Espírito Santo, que ilumina o coração para crer no Invisível, unindo o crente a Cristo e à Igreja como corpo místico.

Na liturgia, especialmente no Sacramento do Batismo, a profissão de fé é central: o catecúmeno renuncia a Satanás e adere ao Deus único, trinitário. Essa recitação do Credo sela o compromisso baptismal, ecoando nas Missas e orações diárias.

A unicidade de Deus, professada contra o politeísmo antigo e moderno, afirma que só Ele é o Senhor absoluto, Criador e Salvador. Essa verdade une judeus e cristãos em uma raiz comum, mas se plenifica na revelação trinitária.

Para adultos a partir dos 20 anos ainda não iniciados nos sacramentos, essa seção convida a uma profissão inicial de fé no RICA, como rito de aceitação, onde expressam o desejo de crer em Deus único, iniciando a jornada de conversão pessoal e comunitária.

\section{Deus Todo-Poderoso (\S 268--278)}

Deus é revelado como \textbf{Pai todo-poderoso}, unindo intimamente sua paternidade amorosa à onipotência absoluta. No Credo, esse atributo é o único nomeado explicitamente entre os divinos, destacando sua relevância para a vida cotidiana: ``De todos os atributos divinos, só a onipotência de Deus é nomeada no Credo: confessar este poder tem grande incidência em nossas vidas.'' (\S 268)

Sua onipotência é \textbf{universal}, abrangendo céu e terra, visível e invisível; governa soberanamente a história e os corações humanos. As Escrituras o proclamam como ``Poderoso de Jacó'', ``Senhor dos exércitos'', que nada impede.

Essa potência se manifesta como \textbf{amor providente}: Deus cuida de suas criaturas como Pai, supre necessidades, perdoa pecados e converte corações rebeldes. ``Pela sua providência ordena todas as coisas convenientemente para a salvação''.

No entanto, a onipotência divina é também \textbf{misteriosa}, atuando na fraqueza --- como na Cruz de Cristo --- e convidando à fé confiante, que nada teme porque tudo está em suas mãos.

Na catequese para adultos não batizados, crer em Deus todo-poderoso dissipa medos e dúvidas modernas, fortalecendo a confiança na providência durante o catecumenato, preparando para receber a força do Espírito no Crisma.

\section{O Criador do Céu e da Terra (\S 279--294)}

A criação é o primeiro testemunho da onipotência divina: Deus chama as coisas do nada ao ser, sem matéria preexistente. ``No princípio, Deus criou o céu e a terra'', revelando sua transcendência como Autor absoluto de tudo. (\S 279)

Essa doutrina rejeita o gnosticismo dualista, que separa Deus do mundo material, e o panteísmo, que o confunde com a criação. Deus é distinto, mas intimamente presente, sustentando o universo a cada instante.

O relato do Gênesis não é mito científico, mas teológico: afirma a bondade originária da criação, ordenada à glória de Deus e ao bem do homem. Cada ser reflete a sabedoria criadora.

A criação continua no tempo pela providência, que guia os eventos para o fim escatológico, combatendo visões fatalistas ou materialistas.

Para catecúmenos adultos, meditar no Criador desperta admiração pelo cosmos, fomentando gratidão e responsabilidade ecológica no RICA, integrando fé e razão em uma visão unitária da realidade.

\section{O Mistério da Criação (\S 295--324)}

Deus criou livremente por amor, para compartilhar sua bem-aventuran\-ça trinitária com criaturas. A criação é um ato de pura bondade, não necessidade, culminando no homem como imagem viva de Deus. (\S 295)

Toda criatura é boa, declarada assim por Deus seis vezes em Gênesis, mas o mal não pertence à ordem criada: surge da liberdade defeituosa, permitido por Deus para um bem maior.

A providência divina governa com sabedoria e ternura, cooperando com a liberdade humana sem violá-la. Deus não abandona sua obra, intervindo milagrosamente quando necessário.

O pecado e o sofrimento não anulam a bondade criadora; ao contrário, a Cruz revela como Deus tira bem do mal, ordenando tudo à redenção.

Em catequese adulta, essa verdade liberta do desespero ante o mal, incentivando oração e obras de misericórdia no período de purificação quaresmal, confiando na providência para a vida cotidiana.

\section{Os Anjos na Criação (\S 325--354)}

Os anjos são criaturas puramente espirituais, inteligentes e livres, criados por Deus para louvá-lo e servi-lo. Numerosos como estrelas, formam o ``céu'' invisível do Credo, mensageiros da vontade divina.

Sua natureza superior à humana os torna guardiões e intercessores, presentes na liturgia e na vida dos santos. Miguel, Gabriel e Rafael exemplificam sua missão salvífica.

Porém, alguns anjos caíram em rebelião soberba, tornando-se demônios, mas sua derrota é certa pelo poder de Cristo.

A doutrina angélica equilibra a visão materialista, afirmando a hierarquia espiritual da criação e a batalha cósmica pelo reino de Deus.

Para adultos em iniciação, conhecer os anjos enriquece a espiritualidade RICA, invocando-os como protetores nos exorcismos menores e na luta contra tentações durante o catecumenato.

\section{O Homem, Imagem de Deus (\S 355--368)}

O homem é coroação da criação: ``Deus criou o homem à sua imagem, à imagem de Deus o criou.'' (\S 355) Une alma espiritual imortal e corpo material, capaz de conhecer, amar e dominar a terra. (\S 356)

Essa imagem confere \textbf{dignidade pessoal} inerente: o homem não é objeto, mas sujeito relacional, dotado de razão, vontade e liberdade. (\S 357) O corpo participa dessa dignidade, expressando a pessoa.

A imortalidade da alma permite comunhão eterna com Deus, enquanto o corpo ressuscitará glorioso. Ferida pelo pecado, a imagem persiste, restaurada pela graça.

A vocação humana é à bem-aventurança, vivendo para Deus em sociedade justa e amorosa.

Na catequese RICA, essa antropologia afirma a dignidade dos adultos em busca de fé, motivando-os a cultivar virtudes no catecumenato para receber plenamente a imagem trinitária nos sacramentos.

\section{Homem e Mulher à Imagem de Deus (\S 369--384)}

Homem e mulher são criados iguais em dignidade, à imagem de Deus em comunhão: ``Homem e mulher os criou.'' (\S 369) Essa reciprocidade reflete a vida trinitária de doação mútua.

O matrimônio originário é imagem do amor esponsal de Cristo pela Igreja, aberto à procriação como participação na criação divina.

A complementaridade sexual não subordina, mas enriquece: cada um realiza-se na alteridade, fundamentando família e sociedade. (\S 371)

A igualdade essencial homem-mulher combate discriminações, promovendo justiça e paz.

Para catecúmenos adultos, essa visão orienta a sexualidade responsável, preparando solteiros ou casados para viver a castidade ou fidelidade conjugal pós-iniciação, integrando corpo e espírito.

\section{A Queda dos Anjos (\S 385--395)}

Deus criou anjos livres; alguns, liderados por Satanás, escolheram soberba, rejeitando submissão ao Criador. ``Tornaram-se demônios por livre escolha, opostos a Deus.'' (\S 391)

Sua ação tenta o homem, mas limitada por Deus, serve à prova da fé e glória maior. A Escritura narra sua derrota na Cruz.

A Igreja ensina vigilância: oração, sacramentos e virtudes combatem o Maligno, sem exageros supersticiosos.

Essa realidade explica tentações universais, sem culpar exclusivamente o diabo pela responsabilidade humana.

No RICA adulto, compreender a queda angélica arma espiritualmente os eleitos, com scrutinios quaresmais libertando de influências malignas para os sacramentos pascais.

\section{O Pecado Original (\S 396--421)}

Adão e Eva, em estado de santidade original, pecaram por desobediência, seduzidos pelo demônio: preferiram autossuficiência a Deus. (\S 396) Romperam a aliança primordial.

Esse pecado, pessoal dos primeiros, transmite-se a todos como pecado original: privação da graça, ferindo natureza com concupiscência, morte e ignorância. (\S 397)

A natureza humana permanece boa, mas inclinada ao mal (\S 398); a Redenção em Cristo a restaura, apagando sua culpa, por meio do Sacramento do Batismo.

Mesmo apagado o pecado original pelo Batismo, suas consequências persistem: a inclinação para o mal exige do homem uma luta interior, necessitando da graça contínua.

Para adultos não iniciados, essa doutrina ilumina a experiência do mal moral no catecumenato, apontando o Batismo como libertação, culminando na mistagogia pascal de nova vida.

\section{Conclusão}
Professar fé no Pai Criador revela o amor originário que sustenta tudo, convidando à conversão total. Adultos em RICA, abracem essa verdade para a iniciação plena.
